\section{\texorpdfstring{Fixed-$u$}{Fixed-u} residue assembly and quantitative continuation}
\label{sec:assembly}

For $\lambda=\alpha+u$, define the complete three-permutation expression
\begin{align}
 \Pcal_\lambda(h,k;t)
 ={}&Z_{\lambda,\beta,\gamma}(h,k)
 +X_{\lambda,\gamma}(t)Z_{-\gamma,\beta,-\lambda}(h,k)
 \notag\\
 &+X_{\beta,\gamma}(t)Z_{\lambda,-\gamma,-\beta}(h,k).
 \label{eq:P-lambda}
\end{align}
The individual terms are meromorphic.  They will always be combined before a
collision divisor is crossed.

The direct and dual diagonal terms for fixed $h,k$ are
\begin{align}
 \mathcal D^+(h,k)
 ={}&\int_{\R}w(t)\frac{1}{2\pi i}\int_{(c_u)}\Gamma(u)Y^u\notag\\
 &\quad\times
 \frac{1}{2\pi i}\int_{(c_s)}\frac{\mathcal G(s,u)}s
 g_{\beta,\gamma}(s,t)
 Z_{\lambda,\beta+s,\gamma+s}(h,k)\dd s\dd u\dd t,
 \label{eq:D-plus}\\
 \mathcal D^-(h,k)
 ={}&\int_{\R}w(t)\frac{1}{2\pi i}\int_{(c_u)}\Gamma(u)Y^u
 X_{\beta,\gamma}(t)\notag\\
 &\quad\times
 \frac{1}{2\pi i}\int_{(c_s)}\frac{\mathcal G(s,u)}s
 g_{-\gamma,-\beta}(s,t)
 Z_{\lambda,-\gamma+s,-\beta+s}(h,k)\dd s\dd u\dd t.
 \label{eq:D-minus}
\end{align}

\subsection{The three residue cancellations}

Reflect the dual diagonal and dual zero mode by $s\mapsto-s$.  Before any
term is separated, the small arithmetic poles at fixed $u$ are
\begin{center}
\begin{tabular}{ll}
\toprule
term & pole locations before reflection where indicated\\
\midrule
$\mathcal D^+$ & $s=-\lambda-\gamma$, $s=-(\beta+\gamma)/2$\\
$\mathcal D^-$ & $s=\beta-\lambda$, $s=(\beta+\gamma)/2$\\
zero mode & $s=-\lambda-\gamma$, $s=\lambda-\beta$\\
\bottomrule
\end{tabular}
\end{center}
The exact identities required at these poles are
\begin{align}
 X_{\beta,\gamma}(t)g_{-\gamma,-\beta}(-s,t)
 &=g_{\beta,\gamma}(s,t)X_{\beta+s,\gamma+s}(t),
 \label{eq:g-diagonal-reflection}\\
 \Phi_{\lambda;\beta,\gamma}(-\lambda-\gamma,t)
 &=g_{\beta,\gamma}(-\lambda-\gamma,t),
 \label{eq:Phi-first-residue}\\
 \Phi_{\lambda;\beta,\gamma}(\lambda-\beta,t)
 &=g_{\beta,\gamma}(\lambda-\beta,t)
 X_{\lambda,\gamma+\lambda-\beta}(t).
 \label{eq:Phi-second-residue}
\end{align}
They follow from \eqref{eq:X-def}, \eqref{eq:Phi-def}, reflection and
duplication.

At $s=-\lambda-\gamma$, the direct diagonal and zero-mode arithmetic factors
both equal
\begin{equation}\label{eq:first-common-Z}
              Z_{\lambda,\beta-\lambda-\gamma,-\lambda}(h,k),
\end{equation}
while their polar zeta arguments have derivatives $+1$ and $-1$.
Equation \eqref{eq:Phi-first-residue} therefore cancels the residues.  After
reflection, the second pair at $s=\lambda-\beta$ has common factor
\begin{equation}\label{eq:second-common-Z}
              Z_{-\gamma-\lambda+\beta,\lambda,-\lambda}(h,k),
\end{equation}
and \eqref{eq:Phi-second-residue} gives opposite residues.  Finally, at
$s_0=-(\beta+\gamma)/2$, the two diagonal factors specialize to
\begin{equation}\label{eq:third-common-Z}
 Z_{\lambda,(\beta-\gamma)/2,(\gamma-\beta)/2}(h,k).
\end{equation}
Their polar arguments have derivatives $+2$ and $-2$, and
$X_{(\beta-\gamma)/2,(\gamma-\beta)/2}(t)=1$.  Thus the third pair also
cancels.  The equalities in \eqref{eq:first-common-Z}--\eqref{eq:third-common-Z}
are equalities of the complete Euler products, not just their global zeta
factors.

Choose
\begin{equation}\label{eq:sigma-s}
 \sigma_s=-\frac12+\eps_s,\qquad
 0<\eps_s<\min\left(\frac14,\frac12-c_u\right),
\end{equation}
and choose $C_*$ so that every signed pair of real shift parts occurring
below has absolute value at most $C_*/\log T$.  For large $T$, take
\begin{equation}\label{eq:cs-geometry}
 c_u+\frac{C_*}{\log T}<c_s<
 \min\left(\frac12-\eps_s,
            1-c_u-\frac{C_*}{\log T}\right).
\end{equation}
In the central strip $-c_s<\Re s<c_s$, the complete reflected integrand has
only
\begin{equation}\label{eq:central-divisors}
 \begin{array}{c|c}
 x=s+\lambda+\gamma=0&s=-\lambda-\gamma,\\
 y=s-\lambda+\beta=0&s=\lambda-\beta,\\
 z=2s+\beta+\gamma=x+y=0&s=-(\beta+\gamma)/2,\\
 s=0&s=0.
 \end{array}
\end{equation}
The $z=1$ locations lie outside by \eqref{eq:cs-geometry} and are zeros of
$\mathcal G$ in any wider deformation.  By \Cref{lem:gamma-collapse}, all
other archimedean poles have imaginary part comparable with $t$ and are
Gaussian-small.

By \Cref{lem:common-contour-recovery}, the four branchwise Tonelli
destinations may now be paired and represented on this common $c_s$-line.
This conclusion was obtained by moving $A$ and the branch line together,
with every small arithmetic divisor kept on the same side; it is not a
branchwise fixed-$A$ contour shift.  Consequently the following contour
difference represents the continuation selected in
\Cref{prop:four-branch-regulator-tonelli}, up to
an error smaller than every power of $T$ after the outer sum.

Put
\begin{equation}\label{eq:F0}
 F_0(s)=\mathcal G(s,u)\Phi_{\lambda;\beta,\gamma}(s,t)
 Z_{-\gamma-s,\beta+s,-\lambda}(h,k).
\end{equation}
Using \eqref{eq:Phi-reflection}, evenness of $\mathcal G$, and symmetry of
$Z$ in its first two shifts, the dual zero mode becomes exactly
\begin{equation}\label{eq:zero-contour-difference}
                  -\frac{1}{2\pi i}\int_{(-c_s)}\frac{F_0(s)}s\dd s.
\end{equation}
The sign comes from $1/(-s)$; reversal of the parametrized contour introduces
no second sign.  Thus the two zero modes form a contour difference, evaluated
by the $x=0$, $y=0$, and $s=0$ residues, with no zero-mode boundary line.

Move the direct diagonal from $(c_s)$ to $(\sigma_s)$ and the reflected dual
diagonal from $(-c_s)$ to $(-\sigma_s)$.  The three nonzero residue pairs
cancel as above.  The $s=0$ residues of the two diagonals and the zero mode
are precisely the three terms in $Q(u)\Pcal_\lambda$.  Undoing the reflection
on the final dual line leaves the full diagonal boundary
\begin{align}
 \mathcal R_{\mathrm{diag},Y}(h,k)
 ={}&\int_{\R}w(t)\frac{1}{2\pi i}\int_{(c_u)}\Gamma(u)Y^u\notag\\
 &\quad\times\bigg\{
 \frac{1}{2\pi i}\int_{(\sigma_s)}\frac{\mathcal G(s,u)}s
 g_{\beta,\gamma}(s,t)\notag\\
 &\hspace{31mm}\times
 Z_{\alpha+u,\beta+s,\gamma+s}(h,k)\dd s\notag\\
 &\quad+X_{\beta,\gamma}(t)
 \frac{1}{2\pi i}\int_{(\sigma_s)}\frac{\mathcal G(s,u)}s
 g_{-\gamma,-\beta}(s,t)\notag\\
 &\hspace{31mm}\times
 Z_{\alpha+u,-\gamma+s,-\beta+s}(h,k)\dd s
 \bigg\}\dd u\dd t.
 \label{eq:diagonal-boundary-Y}
\end{align}
The two terms in \eqref{eq:diagonal-boundary-Y} are separated only after the
joint cancellation at the three small divisors.

Combining \Cref{prop:endpoint-deletion} with this fixed-$u$ residue
calculation yields, on the good patch and after the outer coefficient sum,
\begin{multline}
 \mathcal D^++\mathcal D^-+\mathcal O_{0,Y}^+
 +\mathcal O_{0,Y}^-\\
 =\int_{\R}w(t)\frac{1}{2\pi i}\int_{(c_u)}
 \Gamma(u)Y^uQ(u)\Pcal_{\alpha+u}(h,k;t)\dd u\dd t\\
 {}+\mathcal R_{\mathrm{diag},Y}(h,k)+O_N(T^{-N})
 \label{eq:good-patch-assembly}
\end{multline}
for every $N>0$.  This error includes, in particular, the
coefficient-weighted contribution of
$\mathscr E_{\rm rem}^+(h,k)+\mathscr E_{\rm rem}^-(h,k)$ from
\eqref{eq:zero-plus-evaluated}--\eqref{eq:zero-minus-evaluated}, bounded by
\eqref{eq:remote-remainder-bound}.  As usual in an $h,k$ identity, the error in
\eqref{eq:good-patch-assembly} is understood after summation with
$a_h\overline{b_k}/\sqrt{hk}$.

\paragraph{External-shift objects.}

For external shifts $A,B,C_1$, let $\mathscr A_T(A,B,C_1)$ denote the
complete source: both diagonals and both exact zero modes, with their endpoint
cutoffs, common positive $u$-line, and outer coefficient sum.  Let
\begin{align}
 \mathscr R_T(A,B,C_1)
 ={}&\sum_{h\leq H}\sum_{k\leq K}
 \frac{a_h\overline{b_k}}{\sqrt{hk}}\notag\\
 &\quad\times
 \int_{\R}w(t)\frac{1}{2\pi i}\int_{(c_u)}
 \Gamma(u)Y^uQ(u)\Pcal_{A+u}(h,k;t)\dd u\dd t,
 \label{eq:R-T}\\
 \mathscr D_T(A,B,C_1)
 ={}&\sum_{h\leq H}\sum_{k\leq K}
 \frac{a_h\overline{b_k}}{\sqrt{hk}}
 \mathcal R_{\mathrm{diag},Y}(h,k).
 \label{eq:D-T}
\end{align}
The shifts in every displayed kernel are changed from
$(\alpha,\beta,\gamma)$ to $(A,B,C_1)$.

\begin{lemma}[Rapid-decay propagation]\label{lem:two-constants}
Let $\Omega\subset\C^n$ be a fixed connected domain, let
$U\Subset\Omega$ be nonempty and open, and let $K_0\Subset\Omega$.  If
$f_T$ is holomorphic on $\Omega$ and
\begin{equation}\label{eq:two-constants-hyp}
 \sup_\Omega|f_T|\leq T^{C_2},
 \qquad
 \sup_U|f_T|\ll_NT^{-N}\quad\text{for every }N,
\end{equation}
then, for every $D>0$,
\begin{equation}\label{eq:two-constants-conclusion}
                         \sup_{K_0}|f_T|\ll_{D,K_0}T^{-D}.
\end{equation}
\end{lemma}

\begin{proof}
Join a ball compactly contained in $U$ to $K_0$ by finitely many overlapping
balls in $\Omega$.  Repeated use of the three-ball, or two-constants,
inequality gives
\[
 |f_T|\leq(T^{-N})^{\omega}(T^{C_2})^{1-\omega}
\]
on $K_0$, where $\omega>0$ depends only on the fixed geometry.  Take
$N>(D+(1-\omega)C_2)/\omega$.
\end{proof}

\subsection{Joint removability at shift collisions}

We place the collision algebra before the quantitative continuation that
uses it.  The next lemma is formal local algebra and contains no analytic
continuation input.

\begin{lemma}[Three-divisor cancellation]\label{lem:three-divisor}
Let $x,y$ be independent local holomorphic coordinates and put $z=x+y$.
Suppose $a,b,c,H_0$ are holomorphic and
\begin{equation}\label{eq:three-divisor-form}
 F=\frac{a}{xz}+\frac{b}{yz}+\frac{c}{xy}+H_0,\qquad
 a+c\in(x),\quad b+c\in(y),\quad a-b\in(z).
\end{equation}
Then $F$ is holomorphic across $xyz=0$.
\end{lemma}

\begin{proof}
Over the common denominator, the numerator is
\begin{equation}\label{eq:three-divisor-numerator}
                              N=ay+bx+cz.
\end{equation}
Modulo $x$, use $z=y$ to obtain $N=(a+c)y$; modulo $y$, use $z=x$ to
obtain $N=(b+c)x$; modulo $z$, use $y=-x$ to obtain $N=(b-a)x$.
The elements $x,y,z$ are pairwise nonassociate prime elements in the local
analytic UFD (the corresponding ideals need not be comaximal).  Since each
of these prime elements divides $N$, unique factorization gives
$xyz\mid N$.
\end{proof}

For the fixed-$s$ assembly, take the coordinates in
\eqref{eq:central-divisors}.  After extracting the simple polar zeta factors,
the direct diagonal, reflected dual diagonal, and zero mode have denominator
patterns $xz$, $yz$, and $xy$.  The three residue cancellations in
\eqref{eq:first-common-Z}--\eqref{eq:third-common-Z} give exactly the three
divisibilities in \eqref{eq:three-divisor-form}.  Explicitly,
$\zeta(1+\xi)=\xi^{-1}+h_+(\xi)$ and
$\zeta(1-\xi)=-\xi^{-1}+h_-(\xi)$ with $h_\pm$ holomorphic.  Equality of the
complete Euler residues on $x=0$, $y=0$, and $z=0$ therefore says,
respectively, that the corresponding numerator sums lie in $(x)$, $(y)$,
and $(z)$; these are identities in the local analytic ring, not merely
cancellations at generic points of the three hyperplanes.  Thus the complete fixed-$s$
integrand is jointly holomorphic even at double or triple intersections.  Its
separate $1/s$ factor remains, so the prescribed $s=0$ residue is retained.

For the permutation sum, the collision pairs are
\begin{center}
\begin{tabular}{lll}
\toprule
hyperplane & first Euler product & matching Euler product\\
\midrule
$\lambda=-\gamma$
 &$Z_{-\gamma,\beta,\gamma}$
 &$Z_{-\gamma,\beta,\gamma}$, $X_{-\gamma,\gamma}=1$\\
$\lambda=\beta$
 &$Z_{-\gamma,\beta,-\beta}$
 &$Z_{\beta,-\gamma,-\beta}$, equal by first-shift symmetry\\
$\beta=-\gamma$
 &$Z_{\lambda,-\gamma,\gamma}$
 &$Z_{\lambda,-\gamma,\gamma}$, $X_{-\gamma,\gamma}=1$\\
\bottomrule
\end{tabular}
\end{center}
In each row the polar zeta arguments are opposite local coordinates, so the
residues have opposite signs, including the finite factors at primes dividing
$hk$.

\begin{proposition}\label{prop:P-holomorphic}
The complete expression $\Pcal_\lambda(h,k;t)$ in \eqref{eq:P-lambda} is
jointly holomorphic across all shift-collision divisors.
\end{proposition}

\begin{proof}
Apply \Cref{lem:three-divisor} with
\begin{equation}\label{eq:P-collision-coordinates}
 x=\lambda+\gamma,\qquad
 y=-\beta-\gamma,\qquad
 z=\lambda-\beta=x+y.
\end{equation}
The middle, last, and first terms of \eqref{eq:P-lambda} supply the
denominator patterns $xz$, $yz$, and $xy$, respectively.  The three rows of
the preceding table are precisely the required divisibilities.  This also
handles the common collision locus $\lambda=\beta=-\gamma$, rather than only
generic points of the individual hyperplanes.
\end{proof}

This proposition uses only \eqref{eq:P-lambda}, the Euler-factor symmetry
in \eqref{eq:Z-euler}, and $X_{-\gamma,\gamma}=1$.  In particular, it is now
proved before, and does not import any estimate from, the quantitative
continuation below.

We next justify uniformly the power-factor form quoted after the main
theorem.  Put
\begin{equation}\label{eq:power-X-def}
 p_{A,C_1}(t):=\left(\frac{t}{2\pi}\right)^{-A-C_1}
\end{equation}
and let $\Pcal_\lambda^{\rm pow}$ be \eqref{eq:P-lambda} with
$X_{\lambda,\gamma}$ and $X_{\beta,\gamma}$ replaced by
$p_{\lambda,\gamma}$ and $p_{\beta,\gamma}$.

\begin{proposition}[Collision-safe Stirling replacement]
\label{prop:collision-safe-stirling}
Uniformly in the shift range \eqref{eq:shift-range}, including all pairwise
and triple collisions,
\begin{equation}\label{eq:collision-safe-pointwise}
 \Pcal_\alpha(h,k;t)-\Pcal_\alpha^{\rm pow}(h,k;t)
 \ll_{C,\eps}T^{-1}(\log T)^{C_C}(hk)^\eps
 \qquad(t\in\supp w).
\end{equation}
Consequently its coefficient-weighted $h,k,t$ contribution is
$O_{C,\eps}(T^\eps(HK)^{1/2})$.
\end{proposition}

\begin{proof}
Use the collision coordinates
\[
 x=\lambda+\gamma,\qquad y=-\beta-\gamma,
 \qquad z=\lambda-\beta=x+y,
\]
and write
\[
 \delta_1=X_{\lambda,\gamma}-p_{\lambda,\gamma},
 \qquad
 \delta_2=X_{\beta,\gamma}-p_{\beta,\gamma}.
\]
The complete replacement difference is
\begin{equation}\label{eq:complete-stirling-difference}
 \mathscr D_X
 =\delta_1Z_{-\gamma,\beta,-\lambda}(h,k)
  +\delta_2Z_{\lambda,-\gamma,-\beta}(h,k).
\end{equation}
Set
\[
 A_2=xzZ_{-\gamma,\beta,-\lambda}(h,k),
 \qquad
 A_3=yzZ_{\lambda,-\gamma,-\beta}(h,k).
\]
For $p^\nu\parallel h_0$, put $X_p=p^{-q}$ and $Y_p=p^{-v}$.  The finite
Euler factor in \eqref{eq:Z-euler} is, initially for
$|X_p|,|Y_p|<1$,
\begin{align}
 B_\nu(q,v)
 &:=(1-X_p)(1-Y_p)\sum_{i+j\geq\nu}X_p^iY_p^j \notag\\
 &=1-(1-X_p)(1-Y_p)\sum_{i+j<\nu}X_p^iY_p^j \notag\\
 &=X_p^\nu+(1-X_p)\sum_{j=0}^{\nu-1}X_p^jY_p^{\nu-j}.
 \label{eq:Bnu-symmetric-form}
\end{align}
The middle expression continues $B_\nu$ as a symmetric polynomial in
$X_p,Y_p$.  Consequently \eqref{eq:Z-euler} shows directly that $A_2$ and
$A_3$ are holomorphic near the collision locus.  Since both the exact and
power factors equal $1$ when their two shifts sum to zero,
\[
 \delta_1\in(x),\qquad \delta_2\in(y).
\]
On $z=0$ one has $\lambda=\beta$, hence $\delta_1=\delta_2$; moreover the
finite Euler factors in $A_2$ and $A_3$ agree there by
$B_\nu(q,v)=B_\nu(v,q)$ in \eqref{eq:Bnu-symmetric-form}, equivalently by
the first-shift symmetry in \eqref{eq:Z-euler}.  Since $y=-x$ on this
divisor, the polar factors agree there as the exact identities
\[
 \left.z\zeta(1-z)\right|_{z=0}=-1,
 \qquad
 \left.z\zeta(1+z)\right|_{z=0}=1,
 \qquad
 y\zeta(1+y)=-x\zeta(1-x).
\]
Consequently $A_2=A_3$ on $z=0$.
Thus
\[
 \delta_1A_2\in(x),\qquad
 \delta_2A_3\in(y),\qquad
 \delta_1A_2-\delta_2A_3\in(z).
\]
Applying \Cref{lem:three-divisor} to
\[
 \mathscr D_X=\frac{\delta_1A_2}{xz}
              +\frac{\delta_2A_3}{yz}
\]
proves joint holomorphy, including the triple collision.

It remains to retain the factor $T^{-1}$ uniformly.  Let $L=\log T$ and
choose fixed radii $r_y>r_x>2C$ and $r_\gamma>C$.  The target shift region is
contained in the polydisc
\[
 |x|\leq r_x/L,\qquad |y|\leq r_y/L,\qquad
 |\gamma|\leq r_\gamma/L.
\]
On its distinguished boundary,
$|z|=|x+y|\geq(r_y-r_x)/L$, so no polar divisor is met.  Uniform Stirling
gives $\delta_j\ll_CT^{-1}$ there.  The local formula
\eqref{eq:Z-euler} gives
\[
 Z_{-\gamma,\beta,-\lambda}(h,k),
 Z_{\lambda,-\gamma,-\beta}(h,k)
 \ll_{C,\eps}L^{C_C}(hk)^\eps;
\]
indeed, the final expression in \eqref{eq:Bnu-symmetric-form} shows that,
after the two displayed zeta factors are removed, a prime power
$p^\nu\Vert h_0$ contributes
\[
 \ll_C(\nu+1)p^{-\nu/2+O_C(\nu/L)},
\]
and the exterior $k_0$-power is no larger.  Hence
$\mathscr D_X\ll T^{-1}L^{C_C}(hk)^\eps$ on the distinguished boundary.
The maximum principle in the three variables $(x,y,\gamma)$ transfers this
bound to the whole polydisc and proves
\eqref{eq:collision-safe-pointwise}.  Finally
\[
 \int_\R|w(t)|\dd t\ll T,
 \qquad
 \sum_{h\leq H}\sum_{k\leq K}
 \frac{|a_hb_k|}{\sqrt{hk}}(hk)^\eps
 \ll_\eps T^\eps(HK)^{1/2},
\]
which proves the asserted total error.
\end{proof}

\subsection{Quantitative continuation of the assembled error}

\begin{proposition}[Full-kernel assembly continuation]
\label{prop:full-kernel-continuation}
Under \eqref{eq:main-range}, for every fixed $C_0,D>0$,
\begin{equation}\label{eq:full-continuation-result}
 \mathscr A_T(A,B,C_1)-\mathscr R_T(A,B,C_1)
 -\mathscr D_T(A,B,C_1)=O_D(T^{-D})
\end{equation}
uniformly for
\begin{equation}\label{eq:continuation-polydisc}
                    |A|+|B|+|C_1|\leq C_0/\log T,
\end{equation}
including every shift collision.
\end{proposition}

\begin{proof}
Put $L=\log T$ and use the scaled variables
\begin{equation}\label{eq:scaled-shifts}
                 (\mathsf a,\mathsf b,\mathsf c)=L(A,B,C_1).
\end{equation}
We first fix all geometry.  The admissible contour choices
\begin{equation}\label{eq:majorant-contours}
 c_u=\frac18,\qquad \eps_s=\frac18,\qquad
 \sigma_s=-\frac38,\qquad c_s=\frac14
\end{equation}
satisfy \eqref{eq:cu-range}, \eqref{eq:sigma-s}, and, for sufficiently
large $T$, \eqref{eq:cs-geometry}.  For the prescribed $C_0$, set
\begin{align}
 R_*&=C_0+6,\notag\\
 \Omega&=\{(\mathsf a,\mathsf b,\mathsf c):
       |\mathsf a|,|\mathsf b|,|\mathsf c|<R_*\},
 \label{eq:majorant-domain}\\
 \Omega^\sharp&=\{(\mathsf a,\mathsf b,\mathsf c):
       |\mathsf a|,|\mathsf b|,|\mathsf c|<R_*+1\}.
 \notag
\end{align}
The target compact set
\begin{equation}\label{eq:majorant-target}
 K_0=\{(\mathsf a,\mathsf b,\mathsf c):
       |\mathsf a|+|\mathsf b|+|\mathsf c|\leq C_0\}
\end{equation}
has a fixed collar in $\Omega$.  We use the fixed open good-patch subset
\begin{equation}\label{eq:majorant-good-set}
 U=\left\{ |\mathsf a|<\frac18,
       |\mathsf b-2|<\frac18,
       |\mathsf c+2|<\frac18\right\}.
\end{equation}
Its closure lies in $\Omega$.  The inverse-scaled projection
$(\mathsf b/L,\mathsf c/L)$ of $U$ lies in \eqref{eq:good-patch} with
$\kappa_0=1$ and $c_1=5$, while the
full set $U$ also satisfies the initial three-variable size condition stated
after \eqref{eq:good-patch}, since on $U$ one has
$\Re\mathsf b>15/8$, $\Re\mathsf c<-15/8$, and
$|\mathsf a|+|\mathsf b|+|\mathsf c|<35/8$.

Take $T$ large enough that
\begin{equation}\label{eq:majorant-T-threshold}
                         L>64(R_*+1).
\end{equation}
On $\Omega^\sharp$ the two linear denominator factors in
\eqref{eq:G-two-variable} then have modulus at least $7/8$, and hence their
squares have modulus at least $49/64$.  With $C_*=2(R_*+1)$,
\eqref{eq:majorant-contours} satisfies \eqref{eq:cs-geometry}.  These are
fixed contours throughout $\Omega^\sharp$; the unused kernel divisors stay
a fixed positive distance away.

We now prove a polynomial majorant separately for each assembled object.
In the convention of
\eqref{eq:w-seminorm}--\eqref{eq:auxiliary-weight-seminorm}, all constants
below may depend on
$R_*,B_0,M,\eps$, on the coefficient-bound constants, and on the indicated
finite seminorms, but not on $H,K,T$ or the point of the scaled shift domain.
No new weight-function dependence is introduced in this section.  Choose the
auxiliary divisor-bound exponent sufficiently small.  Once $M$ is fixed, all
factors $L^{C_M}T^{O(\eps)}$ are absorbed into one additional power of $T$.
This last coarse absorption is used only for the polynomial majorant needed
by quantitative continuation, not for the sharp $T^\eps$ in
\eqref{eq:main-asymptotic}.

For the exact-cutoff zero-mode part of $\mathscr A_T$, the mesh bound
\eqref{eq:mesh-majorant} and the full-kernel product bound give
\eqref{eq:source-product-majorant}.  Uniformly on $\Omega^\sharp$,
\begin{align}
 \int_{1/2}^{\infty}l^{-1/2-\Re A}(1+l/Y)^{-P}\dd l
 &\ll_{R_*,B_0}Y^{1/2+(R_*+1)/L}\ll T^{B_0/2},
 \label{eq:majorant-l}\\
 &\sum_{m\geq1}m^{-1/2-\Re B}\bigg\{\notag\\[-2pt]
 &\qquad\int_{1/2}^{\infty}n^{-1/2-\Re C_1}
       (1+mn/T)^{-Q}\dd n\notag\\[-2pt]
 &\qquad{}+(1+m/T)^{-Q}\bigg\}
 \ll_{R_*}T^{1/2}L.
 \label{eq:majorant-mn}
\end{align}
Here $P,Q$ are fixed sufficiently large product-decay orders.  To see the
second estimate, split at $m=T$.  For $m\leq T$ the integral term is at most
\[
 T^{1/2-\Re C_1}
 \sum_{m\leq T}m^{-1-\Re B+\Re C_1}
 \ll_{R_*}T^{1/2}L,
\]
because
$\sum_{m\leq T}m^{-1+2(R_*+1)/L}\ll e^{2(R_*+1)}L$; increasing $P,Q$
gives arbitrary convergence for $m>T$.  This is an exact-cutoff source
bound, not a claim of endpoint-deleted normal convergence in the opposite
half-space.

The outside coefficients cost at most
\begin{equation}\label{eq:majorant-hk}
 \sum_{h\leq H}\sum_{k\leq K}\frac{|a_hb_k|}{\sqrt{hk}}
 \ll_\eps T^\eps(HK)^{1/2}\ll T^{1/2+\eps},
\end{equation}
where $HK\leq T^{1-\eta}$ follows from \eqref{eq:main-range} and $H\geq1$.
Together with the $O(T)$ length of the $t$-integral,
\eqref{eq:majorant-l}--\eqref{eq:majorant-hk} give
$T^{B_0/2+2+o(1)}$ for each exact zero mode.  On the positive line
$\Re s=c_s=1/4$, the two source diagonals are absolutely convergent and the
Euler formula gives the deliberately cruder estimate
\begin{equation}\label{eq:majorant-source-diagonal}
 L^{C_M}T^{1+c_s+B_0c_u}
 H^{1+c_s+\eps}K^{1/2+\eps}\ll T^{B_0/8+4}.
\end{equation}
The exact factors $X_{B,C_1}(t)$ are $O_{R_*}(1)$ on the scaled domain.
The locally uniform exact-cutoff majorants prove holomorphy and
\begin{equation}\label{eq:A-T-polynomial}
                  \sup_{\Omega^\sharp}|\mathscr A_T|
                  \ll T^{B_0/2+4}.
\end{equation}

For $\mathscr R_T$, retain the complete sum $\Pcal_{A+u}$ before taking
absolute values.  On $\Re u=c_u=1/8$, the collision coordinates
\[
 x=A+u+C_1,\qquad y=-B-C_1,\qquad z=A+u-B=x+y
\]
satisfy $|x|,|z|\geq3/32$ on $\Omega^\sharp$ by
\eqref{eq:majorant-T-threshold}.  Thus only $y=0$ can meet this integration
line.  Away from a fixed scaled tube around $y=0$, the Euler formula
\eqref{eq:Z-euler}, Stirling, and vertical-strip zeta bounds give
\begin{equation}\label{eq:P-polynomial}
 \Pcal_{A+u}(h,k;t)
 \ll_{R_*,M,\eps}L^{C_M}(hk)^\eps h^{1/2+\eps}
 (1+|\Im u|)^{C_M}.
\end{equation}
In the scaled tube $|Ly|<1/2$ we first use the logically prior algebraic
input \Cref{prop:P-holomorphic}, with the acyclic dependency recorded above.
More explicitly, fix $A,C_1,u,t$ and put
\[
 f(w)=\left.\Pcal_{A+u}(h,k;t)\right|_{B=-C_1-w/L}.
\]
The complete function $f$ is holomorphic for $|w|<1$ and the circle
$|w|=1/2$ lies inside the fixed collar of $\Omega^\sharp$.  The maximum
principle bounds every $f(w)$ with $|w|\leq1/2$ by the maximum on this
circle.  There
$|B+C_1|=1/(2L)$, so the individual Euler estimates cost only a fixed power
of $L$, while $x$ and $z$ remain bounded away from zero.  A finite cover of
$\overline\Omega$ proves \eqref{eq:P-polynomial} throughout the tube,
including $y=0$, without separately estimating a meromorphic term at the
collision.  The $u,t,h,k$ integrations now give
\begin{equation}\label{eq:R-T-polynomial}
 \sup_\Omega|\mathscr R_T|
 \ll L^{C_M}T^{1+B_0/8}H^{1+\eps}K^{1/2+\eps}
 \ll T^{B_0/8+4}.
\end{equation}

Finally consider $\mathscr D_T$.  This object is defined only after the
direct diagonal, reflected dual diagonal, and zero mode have been assembled
and the three small residues have cancelled.  On
$\Re s=\sigma_s=-3/8$, the two pairs of Euler arguments have real parts
\begin{equation}\label{eq:D-T-Euler-lines}
 \begin{aligned}
 \Re(1+A+u+C_1+s)&=3/4+O(R_*/L),\\
 \Re(1+B+C_1+2s)&=1/4+O(R_*/L),\\
 \Re(1+A+u-B+s)&=3/4+O(R_*/L),\\
 \Re(1-B-C_1+2s)&=1/4+O(R_*/L).
 \end{aligned}
\end{equation}
They stay a fixed distance from both zero and one.  The prime-power formula
\eqref{eq:Bnu-symmetric-form} and vertical-strip zeta bounds therefore give, for both
Euler products on this line,
\begin{equation}\label{eq:D-T-Z-bound}
 Z\ll_{R_*,M,\eps}L^{C_M}(hk)^{1/2+\eps}
 (1+|\Im u|+|\Im s|)^{C_M}.
\end{equation}
Stirling gives
$g(s,t)\ll T^{-3/8}(1+|\Im s|)^{C_M}$; the factor $e^{s^2}$ is Gaussian
and $\Gamma(u)$ has exponential decay.  Completing both vertical integrals
and the $h,k$ sums yields
\begin{equation}\label{eq:D-T-polynomial}
 \sup_\Omega|\mathscr D_T|
 \ll L^{C_M}T^{1-3/8+B_0/8}HK\,T^\eps
 \ll T^{B_0/8+3}.
\end{equation}
This estimate is for the assembled boundary object; no individual
meromorphic beta branch has been continued through a collision.

Define, explicitly,
\begin{equation}\label{eq:Eend-definition}
 \Eend(A,B,C_1)
 :=\mathscr A_T(A,B,C_1)-\mathscr R_T(A,B,C_1)
   -\mathscr D_T(A,B,C_1).
\end{equation}
It follows from \eqref{eq:A-T-polynomial}, \eqref{eq:R-T-polynomial}, and
\eqref{eq:D-T-polynomial} that
\begin{equation}\label{eq:end-error-polynomial}
 F_T(\mathsf a,\mathsf b,\mathsf c)
 :=\Eend(\mathsf a/L,\mathsf b/L,\mathsf c/L),
 \qquad
 \sup_\Omega|F_T|\ll T^{C_2},
 \quad C_2=\frac{B_0}{2}+5.
\end{equation}
Thus the enlarged scaled domain, its collar, and the polynomial exponent are
all fixed independently of $T$.  On $U$, the good-patch assembly
\eqref{eq:good-patch-assembly} gives
$\sup_U|F_T|\ll_NT^{-N}$ for every fixed $N$.

The parameter order is acyclic.  First fix $B_0>2$, $C_0$, and the desired
final saving $D$; this fixes $K_0\Subset\Omega$, $U\Subset\Omega$, and the
propagation exponent $\omega=\omega(C_0)>0$.  Next choose one integer $M$
such that, for $\sigma=M+1/2$,
\begin{equation}\label{eq:M-before-N}
                         (B_0-1)\sigma>D+B_0+20.
\end{equation}
This coarse condition also dominates the later permutation and
diagonal-boundary left-line requirements.  With this $M$ fixed,
$C_2=B_0/2+5$ is fixed; only the lower threshold for $T$ may depend on
$M,C_0$.  Now choose
\begin{equation}\label{eq:N-after-M}
 N=\left\lceil\frac{D+(1-\omega)C_2}{\omega}\right\rceil+2
\end{equation}
and afterward the endpoint integration-by-parts and product-tail orders
needed for the good-patch $O(T^{-N})$.  In particular $M$ is independent of
$N$.  Applying \Cref{lem:two-constants} to $F_T$ on the fixed geometry
$U\Subset\Omega$, with target $K_0$, proves
\eqref{eq:full-continuation-result}.
\end{proof}

Only the difference in \eqref{eq:full-continuation-result} has been
propagated.  In particular, no power-sized estimate for an individual
diagonal boundary term has been continued by unique continuation.
